\section{Introducción teórica}

Este documento es el informe científico-técnico de NeuroMIND, pensado para apoyar el desarrollo diario, la depuración y la trazabilidad metodológica del pipeline de EEG en reposo. El foco analítico se mantiene igual que en las versiones previas del proyecto: conectividad funcional estática a nivel de sensor mediante weighted Phase Lag Index (wPLI), estimada por separado en condiciones de ojos abiertos (EO) y ojos cerrados (EC), y validada mediante inferencia con surrogates de fase aleatorizada.

El objetivo de este cuaderno no es funcionar como manuscrito de publicación, sino como una herramienta de claridad operativa. Por ello, cada sección se organiza alrededor de decisiones de implementación, flujo de datos entre etapas y puntos de control reproducibles. Las descripciones están alineadas con el código en \Julia{} y con los artefactos generados durante las ejecuciones (figuras, tablas, logs y ficheros intermedios serializados), de modo que cualquier fase pueda relanzarse de forma aislada cuando haya que revisar parámetros, controles de calidad o supuestos metodológicos.

El informe conserva la profundidad técnica y la lógica original del pipeline, pero adopta una estructura más limpia para trabajo interno: contexto del proyecto y de los datos, justificación metodológica, especificación fase a fase del pipeline, notas algorítmicas, seguimiento de resultados y cierre operativo con próximos pasos.

\subsection{Esclerosis múltiple: motivación clínica y metodológica}

La esclerosis múltiple (EM) es una enfermedad crónica inflamatoria y desmielinizante del sistema nervioso central en la que procesos inmunomediados alteran la integridad de la mielina, la conducción axonal y el soporte glial, con efectos posteriores sobre la comunicación a gran escala entre vías corticales y subcorticales \citep{Mestre2021}. Desde una perspectiva de sistemas, el impacto no es solo focal: las lesiones y los cambios microestructurales difusos pueden modificar la sincronización, los retrasos de conducción y los patrones de reclutamiento compensatorio.

La motivación de NeuroMIND es precisamente estudiar esa dimensión funcional de red sin asumir que una matriz EEG equivalga a anatomía. En EM se han descrito cambios de conectividad que pueden depender de banda, región, estado de reposo y fase clínica; en algunos contextos se plantea hiperconectividad compensatoria temprana, reorganización interhemisférica o alteración del control top--down, pero estas lecturas requieren muestras comparables, covariables clínicas y control inferencial. En este informe se tratan como hipótesis de trabajo, no como conclusiones de los resultados preliminares.

\begin{figure}[H]
\centering
\includegraphics[width=0.46\textwidth]{theory/cerebro_anatomia_macro.png}
\caption{Anatomía macroscópica como marco visual: lóbulos corticales, cerebelo y tronco encefálico.}
\label{fig:margin-cerebro-anatomia}
\end{figure}

En este contexto, el EEG de superficie se utiliza como una ventana de resolución milisegundo a milisegundo sobre el comportamiento oscilatorio de redes neuronales \citep{Cohen2014,Sanei2022}. El objetivo de este informe no es localizar lesiones, sino caracterizar consecuencias funcionales mediante cambios dependientes de la condición en la organización espectral y el acoplamiento entre canales (EO frente a EC), dentro de un marco reproducible de preprocesamiento e inferencia. Si las escalas EDSS, MSFC u otras variables clínicas quedan disponibles y armonizadas, el siguiente paso natural será relacionarlas con métricas de red por banda de forma exploratoria y posteriormente confirmatoria.

\subsection{Conectividad cerebral}

El EEG en reposo se usa habitualmente para caracterizar dinámicas oscilatorias espontáneas en ausencia de demandas explícitas de tarea, tanto en poblaciones sanas como en contextos patológicos \citep{Cohen2014,Yamada2018}. En este marco, el análisis de conectividad busca cuantificar cómo covarían o se sincronizan las señales entre canales y rangos de frecuencia.

\begin{figure}[H]
\centering
\includegraphics[width=0.46\textwidth]{theory/cerebro_redes_funcionales.png}
\caption{Redes funcionales a gran escala relevantes para interpretar conectividad en EEG.}
\label{fig:margin-cerebro-redes}
\end{figure}

\subsubsection{Niveles conceptuales de conectividad}

La conectividad cerebral puede organizarse en tres niveles complementarios. La conectividad estructural se refiere al cableado anatómico (tractos de sustancia blanca y vías sinápticas) y condiciona qué regiones pueden interactuar. La conectividad funcional cuantifica dependencias estadísticas entre señales, sin imponer explícitamente direccionalidad, y suele estimarse mediante relaciones de covarianza, espectrales o basadas en fase. La conectividad efectiva da un paso más y modela influencia dirigida o causal entre regiones o señales bajo supuestos concretos de modelo.

\begin{fullwidth}
\begin{figure}[H]
\centering
\resizebox{0.86\fulltextwidth}{!}{\NSConnectivityPerspectives}
\caption{Tres lecturas complementarias de una misma red reconstruidas en TikZ: soporte anatómico, sincronía funcional e influencia dirigida.}
\label{fig:connectivity-three-perspectives}
\end{figure}
\end{fullwidth}

Este informe se centra en la \emph{conectividad funcional} a nivel de sensor, porque puede estimarse directamente a partir de EEG en reposo tras un preprocesamiento robusto y control inferencial \citep{Cohen2014,Vinck2011WPLI,Cao2022EEGConnectivity}.

\begin{table}[H]
\centering
\small
\caption{Ejes conceptuales de conectividad y posición actual de NeuroMIND.}
\label{tab:connectivity-concepts}
\begin{tabular}{@{}p{0.22\textwidth}p{0.48\textwidth}p{0.22\textwidth}@{}}
\toprule
\textbf{Eje} & \textbf{Lectura técnica} & \textbf{NeuroMIND actual} \\
\midrule
Estructural & Soporte anatómico: tractos, sinapsis o vías físicas. & No estimada. \\
Funcional & Dependencia estadística entre señales; no implica dirección. & Sí: wPLI por banda. \\
Efectiva & Influencia dirigida/causal bajo un modelo explícito. & No estimada. \\
Dirigida/no dirigida & Aristas con orientación frente a asociaciones simétricas. & No dirigida. \\
Sensor/fuente & Señales en electrodos frente a fuentes reconstruidas. & Sensor; CSD opcional. \\
Estática/dinámica & Una matriz por bloque frente a matrices variables en el tiempo. & Estática por condición. \\
\bottomrule
\end{tabular}
\end{table}

\begin{controlqc}
\small
Las matrices wPLI de NeuroMIND son matrices de conectividad funcional sensor--sensor. No deben interpretarse como conectividad anatómica, dirección de influencia ni causalidad. Una arista alta indica sincronización de fase no nula observada bajo las decisiones de preprocesamiento, segmentación, banda e inferencia, no una vía estructural.
\end{controlqc}

\subsubsection{Formulación en tiempo y frecuencia}

\begin{figure}[H]
\centering
\includegraphics[width=0.46\textwidth]{theory/cerebro_oscilaciones_eeg.png}
\caption{Bandas de oscilación cerebral observables en EEG y lectura funcional general.}
\label{fig:margin-cerebro-oscilaciones}
\end{figure}

Sean $x(t)$ e $y(t)$ dos señales EEG. Su correlación cruzada es
\[
R_{xy}(\tau)=\mathbb{E}[x(t)y(t+\tau)],
\]
con estimador muestral
\[
\hat{R}_{xy}(\tau)=\frac{1}{N}\sum_{k=1}^{N}x_k\,y_{k+\tau}.
\]
En el dominio frecuencial, con $X(f)=\mathcal{F}\{x(t)\}$ e $Y(f)=\mathcal{F}\{y(t)\}$, el espectro cruzado es
\[
S_{xy}(f)=X(f)\,Y^*(f),
\]
donde $(\cdot)^*$ denota conjugación compleja \citep{Cohen2014,Ruth2020,Sanei2022}.

\subsubsection{Métricas no dirigidas de conectividad}

Entre las medidas no dirigidas habituales se incluyen:
\begin{align*}
\mathrm{coh}_{xy}(f) &= \frac{|S_{xy}(f)|}{\sqrt{S_{xx}(f)\,S_{yy}(f)}},\\
\mathrm{PLV} &= \left|\frac{1}{N}\sum_{k=1}^{N} e^{\mathrm{i}\Delta\phi_k}\right|,\\
\mathrm{AEC} &= \mathrm{corr}(A_x(t),A_y(t)).
\end{align*}
La coherencia combina amplitud y fase y es sensible a efectos de desfase cero; PLV se centra en la consistencia de fase; AEC captura cofluctuaciones de envolvente \citep{Cohen2014,Cao2022EEGConnectivity}.

En este proyecto, el estimador principal es wPLI:
\[
\mathrm{wPLI}_{ij}
=
\frac{\left|\mathbb{E}\left[\operatorname{Im}(z_i\overline{z_j})\right]\right|}
{\mathbb{E}\left[\left|\operatorname{Im}(z_i\overline{z_j})\right|\right]+\epsilon},
\]
que reduce el peso de componentes cercanas al desfase cero y encaja mejor con un análisis en espacio de sensores robusto frente a conducción de volumen \citep{Vinck2011WPLI}.

\subsubsection{Justificación de wPLI frente a alternativas}

La elección de wPLI responde a una necesidad concreta del EEG de superficie: estimar sincronización de fase reduciendo el impacto de conducción de volumen y referencia común. Coherencia y PLV pueden aumentar espuriamente cuando dos sensores comparten una fuente o una referencia, porque no distinguen bien acoplamiento fisiológico de contribuciones de desfase cero. La parte imaginaria de la coherencia (ImC) evita ese problema de forma conservadora, pero su magnitud depende del ángulo de fase y puede perder sensibilidad cuando los desfases reales se alejan de su rango óptimo \citep{Vinck2011WPLI}.

PLI mejora esta situación al atender al signo de los desfases no nulos, pero usa una función discontinua: pequeñas perturbaciones alrededor del eje real pueden cambiar una muestra de ``lead'' a ``lag'' y aumentar la sensibilidad al ruido. wPLI conserva la lógica de ignorar desfases compatibles con cero, pero pondera cada contribución por la magnitud de la parte imaginaria del espectro cruzado. Así, las muestras con parte imaginaria pequeña pesan poco, mientras que los desfases no nulos consistentes pesan más \citep{Vinck2011WPLI}.

\begin{ideaclave}
\small
En NeuroMIND, wPLI se usa como compromiso técnico: es menos interpretable que una simple correlación, pero más adecuado para conectividad funcional sensorial cuando se quiere reducir el peso de acoplamientos de fase 0/180 grados. Aun así, wPLI no elimina todos los sesgos: el tamaño muestral, la longitud de segmentos, el número de épocas y la calidad de señal condicionan la estabilidad del estimador.
\end{ideaclave}


\subsubsection{Conectividad efectiva como marco teórico}

Aunque no es el foco principal, los enfoques dirigidos son relevantes como contexto teórico. La causalidad de Granger evalúa si los valores pasados de una señal mejoran la predicción de otra bajo modelos autorregresivos; DTF y PDC cuantifican influencia dirigida en frecuencia dentro de formulaciones MVAR; y DCM plantea el acoplamiento dirigido mediante un modelo generativo explícito. En general, estos métodos de conectividad efectiva dependen de supuestos más fuertes que las métricas funcionales y requieren mayor especificación y validación del modelo \citep{Cohen2014,Sanei2022}.

\subsubsection{Medidas de teoría de la información}

Entre las alternativas basadas en información se encuentran la información mutua y la entropía de transferencia:
\[
I(X;Y)=\sum_{x,y}p(x,y)\log\frac{p(x,y)}{p(x)p(y)},
\]
\[
T_{X\to Y}=\sum p(y_{t+1},y_t,x_t)\log\frac{p(y_{t+1}\mid y_t,x_t)}{p(y_{t+1}\mid y_t)}.
\]
Son medidas flexibles y potencialmente no lineales, pero requieren más datos y mayor coste computacional \citep{Blinowska-Cieslak2022,Cohen2014}.

\subsubsection{Conectividad estática frente a dinámica}

La conectividad estática resume una matriz en ventanas largas:
\[
C_{\text{static}} \in \mathbb{R}^{N\times N}.
\]
La conectividad dinámica estima matrices que varían en el tiempo:
\[
C(t)=f(\text{signal}(t)),\qquad \{C_1,C_2,\dots,C_T\},
\]
habitualmente mediante ventanas deslizantes. En este informe, el análisis principal es estático para ajustarse al diseño de bloques EO/EC y estabilizar las comparaciones inferenciales.

\subsubsection{Por qué wPLI, CSD opcional, surrogates y FDR}

La cadena seleccionada es teóricamente coherente y operativamente trazable: CSD reduce la propagación de desfase cero cuando está activo; wPLI enfatiza la consistencia de fase no nula; los surrogates de fase aleatorizada proporcionan distribuciones nulas empíricas; y FDR controla la multiplicidad por aristas. Por tanto, las matrices resultantes se interpretan como patrones de sincronización estadísticamente controlados y dependientes de condición o grupo, no como vías anatómicas directas ni relaciones causales \citep{Cohen2014,Blinowska-Cieslak2022,Vinck2011WPLI,Friston2011Connectivity,Cao2022EEGConnectivity}.
